\subsection{Principal theta addition and Cartier coefficients in two-rank zero}
\label{subsec:principal-theta-Cartier-bridge}

Throughout this section \(k\) is algebraically closed of characteristic
two.  Let \(B\) be an abelian surface of two-rank zero, let \(M\) be a
symmetric line bundle defining a principal polarization, and choose
\(0\ne\theta\in H^0(B,M)\).  Write
\(F:B\to B^{(2)}\) for relative Frobenius, \(M'=M^{(2)}\),
\(N=M^2=F^*M'\), and \(G=\ker F\).
Fix nonzero invariant volume forms on \(B\) and \(B^{(2)}\).
All Cartier maps below are relative, hence \(k\)-linear; the target
line bundle is always explicitly on the Frobenius twist.

\begin{lemma}[The addition tensor is nondegenerate, including in characteristic two]
\label{lem:principal-theta-addition-perfect}
\status{Proved; independent audits passed.}
Put \(V=H^0(B,N)\), and let
\[
 \xi:B\times B\longrightarrow B\times B,
 \qquad (z,a)\longmapsto(z+a,z-a).
\]
The theorem of the square gives
\(\xi^*(M\boxtimes M)\simeq N\boxtimes N\).
Under any such isomorphism, the nonzero section
\(S=\xi^*(\theta\boxtimes\theta)\) is a tensor in \(V\otimes V\)
whose contraction map \(T_S:V^*\to V\) is an isomorphism.
\end{lemma}
\begin{proof}
The assertion about line bundles follows from the theorem of the cube
and symmetry of \(M\); it is the usual two-variable theorem of the
square, valid over the present field.  The isogeny \(\xi\) has kernel
\(\Delta B[2]\), so pullback from its quotient gives a level subgroup
\[
 s:B[2]\longrightarrow\mathcal G(N\boxtimes N)
\]
above the diagonal.  This is the scheme-theoretic descent/level-subgroup
correspondence of \citet[Theorem~8.10 and Corollary~8.14,
pp.~116--117]{EdixhovenVanDerGeerMoonenAV}.
It fixes \(S\).

The theta group of the external product is the central product
\[
 \mathcal G(N\boxtimes N)
  \simeq(\mathcal G(N)\times\mathcal G(N))
          /\{(c,c^{-1}):c\in\mathbf G_m\}.
\]
For every scheme-valued \(g\in\mathcal G(N)\) over
\(t\in B[2]\), there is a unique \(\tau(g)\in\mathcal G(N)\)
over the same \(t\) such that the class of \((g,\tau(g))\)
is \(s(t)\).  Existence can be checked fppf locally by choosing a lift
over the second factor and adjusting its central scalar; uniqueness
makes the construction descend.  Multiplication of these pairs shows
that \(\tau\) is a homomorphism.  It induces the identity on \(B[2]\)
and inversion on the central \(\mathbf G_m\), hence is an isomorphism.

Let \(\rho\) be the theta action on \(V\).  The tensor \(S\) is
fixed by \(\rho(g)\otimes\rho(\tau(g))\), for all scheme-valued
\(g\).  Consequently \(T_S\) intertwines the dual representation
\(\rho^\vee\) with \(\rho\circ\tau\).  Both are irreducible of
weight minus one.  Indeed, the unrestricted group-scheme theta
representation theorem is Theorem~8.32, not its prime-to-characteristic
predecessor Theorem~8.28; Corollary~9.17 identifies \(H^0(B,N)\) with
its irreducible weight-one representation
\citep[pp.~124--125, 133--134]{EdixhovenVanDerGeerMoonenAV}.
Since \(\xi\) is faithfully flat, \(S\ne0\).  Thus \(T_S\ne0\),
and irreducibility makes it an isomorphism.

We make explicit the group-scheme correction needed in reading the
printed proof of Lemma~8.31 in that source.  The sentence that geometric
\(k\)-characters form a basis of \(\mathcal O(K)\) is not valid for
local-local \(K\).  Instead use the universal character of the perfect
pairing \(K\simeq K^D\): its coefficient functions span
\(\mathcal O(K)\), since the induced linear map
\(\mathcal O(K)^*\to\mathcal O(K)\) is an isomorphism.  Stability
under the diagonal theta-group action, after arbitrary base change,
then makes the subspace appearing in that proof an
\(\mathcal O(K)\)-submodule, hence an ideal after trivializing the
weight-one torsor.  Translation stability leaves only the zero and
whole ideals.  Explicitly, for a Hopf ideal-subspace \(I\) satisfying
\(\Delta(I)\subset I\otimes\mathcal O(K)\), if
\(\varepsilon(I)=0\) the counit identity gives \(I=0\); otherwise
choose \(a\in I\) with \(\varepsilon(a)=1\) and use
\(m(\mathrm{id}\otimes\mathsf S)\Delta(a)=1\) to obtain \(1\in I\).
This is the required universal-character argument; no geometric-point
character spanning claim is imported.
\end{proof}

\begin{lemma}[Surjectivity of the principal Cartier trace]
\label{lem:principal-theta-Cartier-surjective}
\status{Proved; independent audits passed.}
The relative Cartier trace is surjective on these global sections:
\[
 C_M:H^0(B,M^2\otimes\omega_B)
       \longrightarrow H^0(B^{(2)},M'\otimes\omega_{B^{(2)}}).
\]
In particular, after trivializing the invariant volume lines, it is a
nonzero linear functional \(\ell:V\to H^0(B^{(2)},M')\).
\end{lemma}
\begin{proof}
Finite-Frobenius duality identifies the sheaf Cartier map with the
surjective trace \(F_*\omega_B\to\omega_{B^{(2)}}\).  In local
parameters its coefficient formula sends the product of the two
parameters times the volume form to the target volume form; the
finite-duality and relative coefficient conventions are the same as in
the previously checked Cartier-pairing proof
\citep[Proposition~1.3.7 and Theorem~1.3.8,
pp.~26--28]{BrionKumar2005}.
Invariant volumes identify this with a surjection
\(F_*\mathcal O_B\to\mathcal O_{B^{(2)}}\).

The morphism \(F\) is a torsor under \(G\).  Since \(B\) has
two-rank zero, \(G^D\) is connected: it is the kernel of the dual
Verschiebung and has no nonidentity geometric point.  The algebra
\(\mathcal O(G)^*=\mathcal O(G^D)\) is therefore local, and every
simple finite-dimensional \(G\)-representation is trivial.  The
regular representation has a composition series with trivial factors.
Its associated bundle is \(F_*\mathcal O_B\).

The trace kernel has the same kind of filtration.  To see this without
assuming a field trace for the purely inseparable morphism, pull the
surjection back along its torsor \(F\).  Both associated bundles become
trivial, and the pulled-back morphism has constant coefficients because
\(H^0(B,\mathcal O_B)=k\).  Its kernel is consequently the associated
bundle of a subrepresentation of the regular representation.  Its
composition factors are all trivial, so the kernel has a filtration by
\(\mathcal O_{B^{(2)}}\).

Tensor by the ample principal line bundle \(M'\).
The higher cohomology of every resulting factor vanishes, by the
abelian-variety vanishing theorem
\citep[Theorem~9.14, p.~133]{EdixhovenVanDerGeerMoonenAV}.
Thus the trace kernel tensored with \(M'\) has zero first cohomology.
The cohomology exact sequence and projection formula prove the stated
surjectivity.  The target has dimension one.
\end{proof}

\begin{proposition}[The Cartier addition coefficient is a nonzero theta square]
\label{prop:principal-theta-Cartier-square}
\status{Proved; independent audits passed.}
Let \(\theta'\) be the Frobenius twist of \(\theta\).  Applying
relative Cartier in the first variable of the addition tensor gives
\[
 C_z\bigl(S(z,a)\,\omega_z\bigr)
   =c\,\theta'(z)\,\theta(a)^2\,\omega_z',
 \qquad c\in k^*.
\]
The equality is understood through the chosen theorem-of-square
isomorphism and invariant volumes; changing these choices changes the
nonzero scalar.  In particular, its zero divisor as a function of the
translation parameter is exactly \(2\Theta\).
\end{proposition}
\begin{proof}
Cartier in the first variable is a map on \(B^{(2)}\times B\), with
target line bundle \(M'\boxtimes N\).  Since \(h^0(M')=1\), write
its output uniquely as \(\theta'(z)\kappa(a)\omega_z'\), where
\(\kappa\in H^0(B,N)\).  The preceding two results show
\(\kappa\ne0\): it is the contraction of the nondegenerate tensor
\(S\) by the nonzero principal Cartier functional.

It remains to identify its descent.  The section \(S\) is invariant
under the natural diagonal \(G\)-linearization obtained by pullback
through \(\xi\), because \(G\subset B[2]\) and \(\xi\) kills its
diagonal.  On the other hand, \(N=F^*M'\) has its canonical
\(G\)-linearization.  Its external product has a canonical diagonal
\(G\)-linearization on \(N\boxtimes N\).

These two diagonal linearizations agree.  Their difference is an
invertible regular function on \(G\times B\times B\).  Properness,
geometric connectedness, and \(H^0(B\times B,\mathcal O)=k\)
identify this ring of functions with \(\mathcal O(G)\), including
over the nonreduced base \(G\).  The cocycle identity therefore makes
the difference a character \(G\to\mathbf G_m\).  Such a character
is a geometric point of \(G^D\), and there is only the identity point
because \(G^D\) is connected.  Thus \(S\) is invariant for the
canonical diagonal action.

The relative duality trace is unchanged by every automorphism of its
finite-flat cover.  This holds after arbitrary base change, so applies
to every \(T\)-valued \(g\in G(T)\), not only to geometric points.
Translation of the first variable by \(g\) is such an automorphism
of \(F_T\), and the invariant volume form is translation invariant.
Tracing the diagonal invariance of \(S\) consequently shows that
\(\kappa\) is invariant under the canonical \(G\)-linearization
of \(N\).  Faithfully flat descent gives
\[
 H^0(B,N)^G=F^*H^0(B^{(2)},M')=k\,\theta^2.
\]
Since \(\kappa\ne0\), the proportionality scalar is nonzero.
\end{proof}

\begin{corollary}[A conditional odd-cover maximal alternating rank]
\label{cor:principal-theta-odd-cover-dual-rank}
\status{Proved; independent audits passed.}
Let \(\pi:X\to B\) be an odd-degree etale isogeny of degree \(n\)
and put \(L=\pi^*M\).  Write its character decomposition as
\[
 \pi_*\mathcal O_X=\bigoplus_\chi P_\chi,
 \qquad P_\chi P_\psi=P_{\chi\psi}.
\]
For each \(\chi\), let \(a_\chi\in B(k)\) be the unique point
with \(P_\chi\simeq t_{a_\chi}^*M\otimes M^{-1}\).
Suppose \(a_\chi\notin\Theta\) for every nontrivial \(\chi\).
Then some functional on \(H^0(X^{(2)},L^{(2)})\) makes
\[
 (s,t)\longmapsto\lambda\bigl(C_{\rm rel}(st\,\omega_X)\bigr)
\]
an alternating form of rank \(n-1\) on \(H^0(X,L)\).
\end{corollary}
\begin{proof}
Every \(H^0(B,MP_\chi)\) has dimension one.  Choose \(\lambda\)
to vanish on all nontrivial target characters and to be nonzero on the
one-dimensional trivial target summand.  Compatibility of Cartier with
etale pullback and the relative-Frobenius character twist shows that
only the pairs \((\chi,\chi^{-1})\) survive.  The preceding
proposition identifies their coefficients, up to nonzero scalar
trivializations, with \(\theta(a_\chi)^2\), hence they are nonzero.
The trivial summand pairs with itself by zero: the invariant top form
has Cartier zero in two-rank zero, and
\(C(\theta^2\omega)=\theta' C(\omega)\).
There are \((n-1)/2\) inverse pairs, each a nonsingular alternating
two-by-two block.  The resulting rank is \(n-1\).  Alternation for
every section follows also directly from
\(C(s^2\omega_X)=s^{(2)}C(\omega_X)=0\).
\end{proof}

\begin{remark}[The fixed-target direction and the remaining global gate]
\label{rem:principal-theta-Cartier-bridge-scope}
\status{Proved; independent audits passed.}
The theta-avoiding Lagrangian proposition can be applied on \(B\),
then the polarized dual of \(B\to B/H\) supplies an incoming cover
\(X\to B\) whose character points are exactly \(H\).  The preceding
corollary therefore produces maximal dual rank on this auxiliary
\(X\), with \(\pi^*\Theta_B=\ell\Theta_X\).  It does not show that
for a prescribed fixed \(X\) one can choose an outgoing theta-group
quotient whose dual kernel avoids the theta divisor on that quotient.
Nor does maximal dual alternating rank by itself supply the full
global Lang or monodromy argument.  No new fixed-degree full-image
case is asserted here.

The source access for this candidate consists of a complete reread of
the retained author text of AV Theorem~8.10, Corollary~8.14,
Corollary~8.20, the full discussion~8.30--8.32, Theorem~9.14, and
Corollary~9.17 with their proofs.  The local-local character sentence
is explicitly repaired above.  The previously audited finite-Frobenius
duality input is reused; no new original Brion--Kumar access or
mathematical program is claimed.  The full parent proof and source audit passed at 01:27:55.333447 UTC
on 20 November 2026.  The full obstruction/crystal proof and source
audit passed at 01:33:13.137653 UTC.  Both independently checked every
group-scheme, descent, trace and representation step and freshly read
the stated primary sources.  No numerical experiment was required.
The exact candidate and both audit records remain retained.
\end{remark}
